\section{MBR (Master Boot Record)}

Erste 512 Bytes der Festplatte — enthält Bootloader und Partitionstabelle.

\begin{itemize}
	\item \textbf{446 Bytes:} Bootloader-Code
	\item \textbf{4 × 16 Bytes:} Partitionseinträge (maximal 4 primäre Partitionen)
	\item \textbf{2 Bytes:} Signatur \texttt{0x55AA} (Boot-Signatur am Ende)
	\item Maximale Festplattengröße: 2 TB (bzw. 3,2 TB bei 512-Byte-Sektoren)
\end{itemize}

MBR im Hex-Editor ansehen:
\begin{lstlisting}[language=bash]
dd if=/dev/sda count=1 | xxd
\end{lstlisting}

\section{GPT (GUID Partition Table)}

Nachfolger von MBR mit wesentlichen Verbesserungen.

\begin{itemize}
	\item Unterstützt mehr als 4 Partitionen (wichtig für mobile Geräte, z.B. Android)
	\item Keine 2-TB-Grenze
	\item Redundante Kopie der Partitionstabelle am Ende der Festplatte
	\item Enthält Protective MBR für Abwärtskompatibilität
\end{itemize}

\section{mmls — Partitionstabelle interpretieren}

\texttt{mmls} zeigt die Partitionstabelle eines Images.
Die Ausgabe liefert den \textbf{Offset in Sektoren}, der für alle TSK-Tools als \texttt{-o} Parameter benötigt wird.

\begin{lstlisting}[language=bash]
mmls image.E01
\end{lstlisting}

Beispielausgabe:
\begin{lstlisting}
      Slot   Start        End          Length       Description
000:  Meta   0000000000   0000000000   0000000001   Primary Table (#0)
001:  ----   0000000001   0000002047   0000002047   Unallocated
002:  000    0000002048   0000206847   0000204800   EFI system partition
003:  001    0000206848   0248655871   0248449024   Basic data partition
\end{lstlisting}

\begin{description}
	\item[Start] Offset in Sektoren — diesen Wert als \texttt{-o} bei TSK-Tools verwenden
	\item[Length] Länge in Sektoren — \texttt{Length × 512} = Partitionsgröße in Bytes
	\item[DOS] Partitionstabellen-Typ ist MBR
	\item[GPT] Partitionstabellen-Typ ist GPT
	\item[Meta] Metadaten-Einträge (Tabelle selbst)
	\item[----] Unallocated Space zwischen Partitionen
\end{description}

Beispiel: Dateisystem der dritten Partition analysieren:
\begin{lstlisting}[language=bash]
fsstat -o 206848 image.E01
fls -pr -o 206848 image.E01
\end{lstlisting}
